\documentclass[11pt]{article}
\usepackage{amsmath,amssymb,amsthm}
\usepackage{graphicx}
\usepackage[utf8]{inputenc}
\usepackage{geometry}
\geometry{a4paper, margin=1in}

\theoremstyle{definition}
\newtheorem{definition}{Definition}[section]
\newtheorem{example}[definition]{Example}
\theoremstyle{plain}
\newtheorem{theorem}[definition]{Theorem}
\newtheorem{lemma}[definition]{Lemma}
\newtheorem{proposition}[definition]{Proposition}
\newtheorem{corollary}[definition]{Corollary}

\newcommand{\stag}[1]{}
\newenvironment{construction}{}{}

\title{Homework: Group Structures and Morphisms}
\author{Algebra II}
\date{\today}

\begin{document}

\section{Group structure and morphisms}

% An included PDF figure. Visual Mode has to show the rendered page here, not a
% placeholder, so the smoke probe asserts that this canvas has real pixels in it.
% `figure.pdf` is generated by `scripts/make-fixture-pdf.mjs`.
\begin{figure}[h]
\centering
\includegraphics[width=0.35\textwidth]{figure.pdf}
\caption{A commuting square, drawn as a PDF.}\label{fig:square}
\end{figure}

\begin{definition}\label{grp:group}
A group is a set $G$ with a multiplication $G \times G \to G$, an element $1$, and an inverse operation $g \mapsto g^{-1}$ satisfying
\[
  (gh)k = g(hk), \quad 1g = g = g1, \quad gg^{-1} = 1 = g^{-1}g.
\]
It is abelian if $gh = hg$ for all $g, h \in G$. For abelian groups we often write the operation as addition, the identity as $0$, and the inverse of $a$ as $-a$. The one-element group is the trivial group.
\end{definition}

\begin{lemma}[Elementary group calculus]\label{grp:calculus}
In a group the identity and each inverse are unique. Left and right cancellation hold, $(gh)^{-1} = h^{-1}g^{-1}$, and $(g^{-1})^{-1} = g$. The equations $ax = b$ and $ya = b$ have the unique solutions $x = a^{-1}b$ and $y = ba^{-1}$. Define $g^0 = 1$, $g^{n+1} = g^n g$ for $n \ge 0$, and $g^{-n} = (g^{-1})^n$. For integers $m, n$, $g^{m+n} = g^m g^n$ and $(g^m)^n = g^{mn}$. If $gh = hg$, then $(gh)^n = g^n h^n$ for every integer $n$.
\end{lemma}

\begin{example}[Elementary groups]\label{grp:elementary-examples}
The integers, rational numbers, and real numbers under addition are abelian groups. For an integer $n \ge 1$, addition of residue classes makes $\mathbb{Z}/n\mathbb{Z}$ an abelian group of order $n$; write $C_n$ for its isomorphism type. The nonzero rational numbers under multiplication form an abelian group. The natural numbers under addition are not a group, since positive elements have no additive inverses there.
\end{example}

\begin{definition}\label{grp:homomorphism}
A homomorphism $f : G \to H$ satisfies $f(gg') = f(g)f(g')$. An isomorphism is a homomorphism with a two-sided inverse homomorphism. An endomorphism is a homomorphism $G \to G$; an automorphism is an isomorphism $G \to G$. Write $\mathrm{Hom}(G, H)$, $\mathrm{End}(G)$, and $\mathrm{Aut}(G)$ for these sets.
\end{definition}

\section{Cosets and Normal Subgroups}

Let $H$ be a subgroup of $G$. The left cosets of $H$ in $G$ partition the group $G$.
Each coset $gH$ has cardinality equal to $|H|$, leading to Lagrange's theorem.
When $gH = Hg$ for every $g \in G$, $H$ is normal, denoted $H \trianglelefteq G$.

\subsection{Quotient Group Construction}

For $N \trianglelefteq G$, the cosets $G/N$ form a group under the operation
$(aN)(bN) = (ab)N$. The canonical projection $\pi : G \to G/N$ given by $\pi(g) = gN$
is a surjective group homomorphism with $\ker \pi = N$.

The quotient construction is characterised by the following properties:

\begin{itemize}
  \item The projection $\pi$ is surjective, and $\pi(g) = \pi(g')$ exactly when
        $g^{-1}g' \in N$.
  \item $N$ is the kernel of $\pi$, so $N$ is normal in $G$.
  \item Every homomorphism $f : G \to H$ with $N \le \ker f$ factors uniquely
        through $\pi$.
\end{itemize}

Equivalently, one may check the three conditions below in order:

\begin{enumerate}
  \item $N$ is a subgroup of $G$.
  \item $gN = Ng$ for every $g \in G$.
  \item The product $(aN)(bN) = (ab)N$ is well defined.
\end{enumerate}

% Mathematical commentary, exercises and background notes on quotient structures (1)
% Mathematical commentary, exercises and background notes on quotient structures (2)
% Mathematical commentary, exercises and background notes on quotient structures (3)
% Mathematical commentary, exercises and background notes on quotient structures (4)
% Mathematical commentary, exercises and background notes on quotient structures (5)
% Mathematical commentary, exercises and background notes on quotient structures (6)
% Mathematical commentary, exercises and background notes on quotient structures (7)
% Mathematical commentary, exercises and background notes on quotient structures (8)
% Mathematical commentary, exercises and background notes on quotient structures (9)
% Mathematical commentary, exercises and background notes on quotient structures (10)
% Mathematical commentary, exercises and background notes on quotient structures (11)
% Mathematical commentary, exercises and background notes on quotient structures (12)
% Mathematical commentary, exercises and background notes on quotient structures (13)
% Mathematical commentary, exercises and background notes on quotient structures (14)
% Mathematical commentary, exercises and background notes on quotient structures (15)
% Mathematical commentary, exercises and background notes on quotient structures (16)
% Mathematical commentary, exercises and background notes on quotient structures (17)
% Mathematical commentary, exercises and background notes on quotient structures (18)
% Mathematical commentary, exercises and background notes on quotient structures (19)
% Mathematical commentary, exercises and background notes on quotient structures (20)
% Mathematical commentary, exercises and background notes on quotient structures (21)
% Mathematical commentary, exercises and background notes on quotient structures (22)
% Mathematical commentary, exercises and background notes on quotient structures (23)
% Mathematical commentary, exercises and background notes on quotient structures (24)
% Mathematical commentary, exercises and background notes on quotient structures (25)
% Mathematical commentary, exercises and background notes on quotient structures (26)
% Mathematical commentary, exercises and background notes on quotient structures (27)
% Mathematical commentary, exercises and background notes on quotient structures (28)
% Mathematical commentary, exercises and background notes on quotient structures (29)
% Mathematical commentary, exercises and background notes on quotient structures (30)
% Mathematical commentary, exercises and background notes on quotient structures (31)
% Mathematical commentary, exercises and background notes on quotient structures (32)
% Mathematical commentary, exercises and background notes on quotient structures (33)
% Mathematical commentary, exercises and background notes on quotient structures (34)
% Mathematical commentary, exercises and background notes on quotient structures (35)
% Mathematical commentary, exercises and background notes on quotient structures (36)
% Mathematical commentary, exercises and background notes on quotient structures (37)
% Mathematical commentary, exercises and background notes on quotient structures (38)
% Mathematical commentary, exercises and background notes on quotient structures (39)
% Mathematical commentary, exercises and background notes on quotient structures (40)
% Mathematical commentary, exercises and background notes on quotient structures (41)
% Mathematical commentary, exercises and background notes on quotient structures (42)
% Mathematical commentary, exercises and background notes on quotient structures (43)
% Mathematical commentary, exercises and background notes on quotient structures (44)
% Mathematical commentary, exercises and background notes on quotient structures (45)
% Mathematical commentary, exercises and background notes on quotient structures (46)
% Mathematical commentary, exercises and background notes on quotient structures (47)
% Mathematical commentary, exercises and background notes on quotient structures (48)
% Mathematical commentary, exercises and background notes on quotient structures (49)
% Mathematical commentary, exercises and background notes on quotient structures (50)
% Mathematical commentary, exercises and background notes on quotient structures (51)
% Mathematical commentary, exercises and background notes on quotient structures (52)
% Mathematical commentary, exercises and background notes on quotient structures (53)
% Mathematical commentary, exercises and background notes on quotient structures (54)
% Mathematical commentary, exercises and background notes on quotient structures (55)
% Mathematical commentary, exercises and background notes on quotient structures (56)
% Mathematical commentary, exercises and background notes on quotient structures (57)
% Mathematical commentary, exercises and background notes on quotient structures (58)
% Mathematical commentary, exercises and background notes on quotient structures (59)
% Mathematical commentary, exercises and background notes on quotient structures (60)
% Mathematical commentary, exercises and background notes on quotient structures (61)
% Mathematical commentary, exercises and background notes on quotient structures (62)
% Mathematical commentary, exercises and background notes on quotient structures (63)
% Mathematical commentary, exercises and background notes on quotient structures (64)
% Mathematical commentary, exercises and background notes on quotient structures (65)
% Mathematical commentary, exercises and background notes on quotient structures (66)
% Mathematical commentary, exercises and background notes on quotient structures (67)
% Mathematical commentary, exercises and background notes on quotient structures (68)
% Mathematical commentary, exercises and background notes on quotient structures (69)
% Mathematical commentary, exercises and background notes on quotient structures (70)
% Mathematical commentary, exercises and background notes on quotient structures (71)
% Mathematical commentary, exercises and background notes on quotient structures (72)
% Mathematical commentary, exercises and background notes on quotient structures (73)
% Mathematical commentary, exercises and background notes on quotient structures (74)
% Mathematical commentary, exercises and background notes on quotient structures (75)
% Mathematical commentary, exercises and background notes on quotient structures (76)
% Mathematical commentary, exercises and background notes on quotient structures (77)
% Mathematical commentary, exercises and background notes on quotient structures (78)
% Mathematical commentary, exercises and background notes on quotient structures (79)
% Mathematical commentary, exercises and background notes on quotient structures (80)
% Mathematical commentary, exercises and background notes on quotient structures (81)
% Mathematical commentary, exercises and background notes on quotient structures (82)
% Mathematical commentary, exercises and background notes on quotient structures (83)
% Mathematical commentary, exercises and background notes on quotient structures (84)
% Mathematical commentary, exercises and background notes on quotient structures (85)
% Mathematical commentary, exercises and background notes on quotient structures (86)
% Mathematical commentary, exercises and background notes on quotient structures (87)
% Mathematical commentary, exercises and background notes on quotient structures (88)
% Mathematical commentary, exercises and background notes on quotient structures (89)
% Mathematical commentary, exercises and background notes on quotient structures (90)
% Mathematical commentary, exercises and background notes on quotient structures (91)
% Mathematical commentary, exercises and background notes on quotient structures (92)
% Mathematical commentary, exercises and background notes on quotient structures (93)
% Mathematical commentary, exercises and background notes on quotient structures (94)
% Mathematical commentary, exercises and background notes on quotient structures (95)
% Mathematical commentary, exercises and background notes on quotient structures (96)
% Mathematical commentary, exercises and background notes on quotient structures (97)
% Mathematical commentary, exercises and background notes on quotient structures (98)
% Mathematical commentary, exercises and background notes on quotient structures (99)
% Mathematical commentary, exercises and background notes on quotient structures (100)
% Mathematical commentary, exercises and background notes on quotient structures (101)
% Mathematical commentary, exercises and background notes on quotient structures (102)
% Mathematical commentary, exercises and background notes on quotient structures (103)
% Mathematical commentary, exercises and background notes on quotient structures (104)
% Mathematical commentary, exercises and background notes on quotient structures (105)
% Mathematical commentary, exercises and background notes on quotient structures (106)
% Mathematical commentary, exercises and background notes on quotient structures (107)
% Mathematical commentary, exercises and background notes on quotient structures (108)
% Mathematical commentary, exercises and background notes on quotient structures (109)
% Mathematical commentary, exercises and background notes on quotient structures (110)
% Mathematical commentary, exercises and background notes on quotient structures (111)
% Mathematical commentary, exercises and background notes on quotient structures (112)
% Mathematical commentary, exercises and background notes on quotient structures (113)
% Mathematical commentary, exercises and background notes on quotient structures (114)
% Mathematical commentary, exercises and background notes on quotient structures (115)
% Mathematical commentary, exercises and background notes on quotient structures (116)
% Mathematical commentary, exercises and background notes on quotient structures (117)
% Mathematical commentary, exercises and background notes on quotient structures (118)
% Mathematical commentary, exercises and background notes on quotient structures (119)
% Mathematical commentary, exercises and background notes on quotient structures (120)
% Mathematical commentary, exercises and background notes on quotient structures (121)
% Mathematical commentary, exercises and background notes on quotient structures (122)
% Mathematical commentary, exercises and background notes on quotient structures (123)
% Mathematical commentary, exercises and background notes on quotient structures (124)
% Mathematical commentary, exercises and background notes on quotient structures (125)
% Mathematical commentary, exercises and background notes on quotient structures (126)
% Mathematical commentary, exercises and background notes on quotient structures (127)
% Mathematical commentary, exercises and background notes on quotient structures (128)
% Mathematical commentary, exercises and background notes on quotient structures (129)
% Mathematical commentary, exercises and background notes on quotient structures (130)
% Mathematical commentary, exercises and background notes on quotient structures (131)
identities and composition, so quotient formation is covariant in such
maps of pairs $(G, N)$.
\end{construction}
\begin{theorem}[First Isomorphism Theorem]
\stag{grp:first-isomorphism}
For every homomorphism $f : G \to H$, the map
$G/\ker f \to f(G)$, $g\ker f \mapsto f(g)$, is an isomorphism.
Thus $f$ factors as a quotient, an isomorphism, and a subgroup inclusion.
The kernel inclusion also has a mapping property: every homomorphism
$u : K \to G$ with $fu$ trivial factors uniquely through $\ker f$.
\end{theorem}
\begin{theorem}[Subgroup Correspondence]
\stag{grp:correspondence}
For $N \trianglelefteq G$ and quotient map $q$, the assignments
$H \mapsto H/N$ and $L \mapsto q^{-1}(L)$ are mutually inverse,
inclusion-preserving bijections between subgroups of $G$ containing
$N$ and subgroups of $G/N$. They preserve intersections and generated
joins and restrict to a correspondence of normal subgroups. For
$N \le H \le G$, $[G : H] = [G/N : H/N]$.
\end{theorem}
\begin{theorem}[Second And Third Isomorphism Theorems]
\stag{grp:other-isomorphisms}
If $H \le G$ and $N \trianglelefteq G$, then $HN = NH$ is a subgroup,
$N \trianglelefteq HN$, $H \cap N \trianglelefteq H$, and
\[
  H/(H \cap N) \xrightarrow{\cong} HN/N, \quad h(H \cap N) \mapsto hN.
\]
\end{theorem}

\end{document}
